\documentclass[11pt]{article}
\usepackage[review]{acl}

\usepackage{times}
\usepackage{latexsym}
\usepackage[T1]{fontenc}
\usepackage[utf8]{inputenc}
\usepackage{microtype}
\usepackage{inconsolata}
\usepackage{graphicx}
\usepackage{booktabs}
\usepackage{amsmath}
\usepackage{amssymb}
\usepackage{amsthm}
\usepackage{enumitem}
\usepackage{multirow}
\usepackage{xspace}
\usepackage{array}
\usepackage{tikz}
\usetikzlibrary{shapes,arrows,positioning,calc}
\usepackage{pgfplots}
\pgfplotsset{compat=1.18}

\newcommand{\system}{\textsc{EVIRAG}\xspace}
\newcommand{\bench}{\textsc{EVIRAG-Bench}\xspace}
\setlength\titlebox{6cm}

\newtheorem{definition}{Definition}
\newtheorem{remark}{Remark}

%% ─────────────────────────────────────────────────────────────────────────────
\title{Against Epistemic Collapse: Disagreement-Aware Scientific\\
       Retrieval-Augmented Generation}

\author{Anonymous Submission}

\begin{document}
\maketitle

%% ─────────────────────────────────────────────────────────────────────────────
\begin{abstract}
Retrieval-augmented generation (RAG) has become the default interface between
users and scientific corpora --- yet every standard RAG pipeline makes a hidden
architectural bet: that scientific evidence converges on one answer.  When the
evidence is genuinely contested, this bet produces what we call \emph{epistemic
collapse}: the systematic conversion of structured multi-view disagreement into
a single confident response.  We show, both analytically and empirically, that
epistemic collapse is structurally unavoidable for single-view systems: any
output constrained to one embedding cannot simultaneously cover mutually
distinct viewpoints.

We introduce \system\ (Evidence-centric, Viewpoint-Inclusive RAG), a RAG
framework that treats scientific controversy as a \emph{primary output
signal} rather than noise to suppress.  \system\ combines adversarial
multi-agent retrieval, claim-level disagreement-graph construction, a novel
causal disagreement attribution taxonomy (CDA-7), temporal epistemic drift
tracking, multi-view synthesis, and evidence-structured confidence calibration.
We contribute: (i)~a four-metric formal suite (CCS, ED, PI, EC@K) with
formal metric definitions and empirically validated boundary properties,
(ii)~a four-class controversy taxonomy derived from the ED--PI space,
and (iii)~\bench, an evaluation benchmark for epistemic fidelity in
scientific QA comprising 250 expert-annotated queries across five domains.  Across a 250-query evaluation spanning education, biomedicine,
economics, earth sciences, and nutrition with \texttt{qwen3.6:35b-a3b},
\system\ achieves CR$\,{=}\,0.742$,
a $+0.170$ improvement over vanilla RAG
($p{<}0.001$, Wilcoxon signed-rank, Bonferroni-corrected),
$+0.124$ over MADAM-RAG ($p{=}0.003$),
and $+0.189$ over PaperQA2 ($p{<}0.001$);
dense-encoder ViewpointCoverage rises from $0.423$ (vanilla) to $0.847$;
CCS improves from $0.571$ to $0.518$ ($p{=}0.008$);
FaithfulnessScore reaches $0.891$.
Human evaluation (50 instances, 3 expert annotators, Krippendorff's
$\alpha{=}0.72$) confirms: \system\ Full scores $4.3/5$ vs.\ vanilla
RAG's $2.1/5$ on epistemic completeness.  A fast offline claim-graph
mode provides $25.62{\times}$ latency improvement.
\end{abstract}

%% ─────────────────────────────────────────────────────────────────────────────
\section{Introduction}
\label{sec:intro}

\paragraph{The problem is structural, not incidental.}
Consider three scientific questions across domains: \textit{Does homework
improve academic achievement?} \textit{Do statins prevent cardiovascular events
in primary prevention?} \textit{Does immigration raise native unemployment?}
All three have large empirical literatures.  All three are \emph{genuinely
contested}: experts disagree not because of missing data but because of
systematic methodological, population, and theoretical differences.  A standard
RAG pipeline, applied to any of these, returns one fluent answer.

That answer is locally supported.  It is also epistemically incomplete in a way
that cannot be fixed by improving retrieval recall or generation quality.
If the output interface permits only one answer, disagreement is erased
\emph{by construction}.

\paragraph{Concretely.}
On the homework question, a standard pipeline returns: ``Research consistently
shows that homework improves academic performance.''  \system\ returns:

\begin{quote}
\small
\textbf{Dominant view} [5 claims, 3 sources]: \emph{``Homework produces
conditional academic benefits dependent on grade level ($\geq$7), subject, and
assignment design.''}
\textbf{Alternative view} [3 claims, 2 sources]: \emph{``No conclusive
achievement gains attributable to homework across the meta-analytic record.''}
\textbf{Minority view} [2 claims, 2 sources]: \emph{``Excess homework causes
measurable wellbeing harm without commensurate academic gain.''}
Disagreement density:~5.3\%.  Confidence: \textsc{medium}.
Controversy class: \textit{stable} (ED$\!=\!0.38$, PI$\!=\!0.12$).
CDA-7 attribution: dominant$\!\leftrightarrow\!$alternative is
\textsc{methodological}; minority divergence is \textsc{operational}.
\end{quote}

The single-view answer covers 1 of 3 viewpoints: ViewpointCoverage $= 1/3$.
Intuitively, any system constrained to a single output embedding cannot
simultaneously maximize similarity to mutually orthogonal viewpoints --- the
output must compromise, systematically underrepresenting minority views
(\S\ref{sec:formalization}).

\paragraph{Why existing approaches do not solve this.}
Multi-perspective summarization (e.g., \citealt{bravzinskas2020few,
hayashi2021wikiasp}) generates diverse summaries from opinionated text (reviews,
arguments), but lacks per-view provenance, source-level attribution, causal
disagreement typing, or temporal tracking --- capabilities essential for
scientific corpora where \emph{why} papers disagree determines how to communicate
uncertainty.  Scientific claim verification (SciFact, PaperQA2) adjudicates
individual claims against evidence; adjudication presupposes that controversy is
resolvable.  Conflicting-evidence RAG (MADAM-RAG) resolves conflict by
multi-agent debate; we argue that for genuine scientific controversy, resolution
is premature.  None of these systems surfaces the \emph{structure} of
disagreement as a first-class output.

\paragraph{Our approach.}
We study epistemic fidelity in scientific RAG: whether systems preserve
structured disagreement rather than collapsing contested evidence into a single
answer.  \system\ enforces epistemic fidelity by (i)~explicitly breaking the
consensus attractor via adversarial multi-agent retrieval, (ii)~structuring
evidence as a disagreement graph before any synthesis, (iii)~attributing the
causes of each contradiction via CDA-7, (iv)~tracking temporal trajectory of
controversy, and (v)~synthesizing a multi-view answer surface with confidence
calibrated to evidence structure, not output fluency.

\paragraph{Contributions.}
\begin{enumerate}[leftmargin=*,nosep]
  \item \textbf{Epistemic collapse formalization (\S\ref{sec:formalization}):}
    We define epistemic collapse (Definition~\ref{def:collapse}) and show
    analytically and empirically that it is structurally unavoidable for
    single-view systems and amplified by topical-relevance retrieval.
    These structural arguments motivate the \system\ architecture and are
    directly confirmed by the empirical results (\S\ref{sec:main-results}).
  \item \textbf{Four-metric formal suite (\S\ref{sec:metrics}):} Consensus
    Collapse Score (CCS), Epistemic Divergence (ED), Polarization Index (PI),
    and EpistemicCoverage@K (EC@K) --- a metric suite designed to
    measure information loss from viewpoint compression, with Fiedler-inspired
    spectral grounding for PI and formally characterized boundary properties
    for CCS (validated against human epistemic completeness at $\rho{=}{-}0.82$).
  \item \textbf{The \system\ framework (\S\ref{sec:system}):} A seven-stage
    architecture with adversarial retrieval, claim-graph reasoning, CDA-7
    causal attribution, temporal drift tracking, multi-view synthesis, and
    evidence-structured confidence.
  \item \textbf{CDA-7 taxonomy (\S\ref{sec:cda7}):} A seven-class causal
    disagreement attribution taxonomy grounded in science and technology studies
    literature \citep{longino2002science,solomon2001social}, with resolvability
    priors that modulate confidence calibration.  To our knowledge, no prior RAG
    pipeline operationalizes causal attribution typing at this level.
  \item \textbf{\bench\ and evaluation (\S\ref{sec:eval}):} We contribute an
    evaluation benchmark for epistemic fidelity in scientific QA, comprising
    250 expert-annotated queries across five domains, evaluated against
    three external baselines (MADAM-RAG, PaperQA2, MMR-RAG) and four
    internal ablation variants, with human evaluation ($n{=}750$ ratings,
    Krippendorff's $\alpha{=}0.72$) and a user trust calibration
    experiment.  A fast offline claim-graph mode achieves $25.62{\times}$
    latency reduction over vanilla RAG.
\end{enumerate}

%% ─────────────────────────────────────────────────────────────────────────────
\section{Related Work}
\label{sec:related}

\paragraph{RAG and attribution.}
\citet{lewis2020rag} established retrieve-then-generate.  RARR improves
post-hoc attribution \citep{gao2023rarr}; FActScore evaluates claim-level
factual precision \citep{min2023factscore}; RAGChecker provides diagnostic
decomposition \citep{ru2024ragchecker}.  All retain a single-view output
interface and treat attribution as a quality check on one answer.

\paragraph{Scientific claim verification.}
SciFact and SciFact-Open cast scientific reasoning as binary claim verification
\citep{wadden2020scifact,wadden2022scifactopen}.  ExpertQA curates expert
questions with attributed answers \citep{malaviya2024expertqa}.  PaperQA2
applies agentic chain-of-thought retrieval to scientific questions
\citep{lala2023paperqa}.  All systems adjudicate or verify: a single answer is
produced.  Adjudication presupposes that the controversy is resolvable --- the
exact assumption \system\ rejects.

\paragraph{Conflicting-evidence and multi-answer QA.}
RoMQA studies multiple correct answers from multiple evidence sources
\citep{zhong2023romqa}.  MADAM-RAG uses multi-agent debate to select the best
surviving answer under conflict \citep{wang2025conflicting}.  These systems
\emph{resolve} conflict; we \emph{preserve} it.

\paragraph{Multi-perspective summarization.}
A parallel line of work generates multiple summaries from opinionated corpora:
\citet{bravzinskas2020few} synthesize few-shot multi-perspective summaries from
product reviews; \citet{hayashi2021wikiasp} produce aspect-based Wikipedia
summaries from multiple sources.  These systems produce diverse output text but
are designed for review/argument corpora rather than scientific literature,
lack CDA-7-style causal attribution of \emph{why} perspectives differ, do not
model temporal evolution of disagreement, and provide no evidence-structured
confidence calibration.  Viewpoint diversity in multi-perspective summarization
is a generation objective; in \system\ it is a retrieval-and-structure objective
grounded in formal metrics.

\paragraph{Epistemic uncertainty in LLMs.}
\citet{xiong2024llms} and \citet{kadavath2022language} study how confidently
LLMs express individual statements.  These approaches address the generation
model's internal uncertainty.  \system\ operates at a different layer: the
epistemic structure of the \emph{retrieved evidence corpus}, which is
orthogonal to model calibration.

\paragraph{Philosophy of scientific disagreement.}
The CDA-7 taxonomy is grounded in science and technology studies work on the
causes of scientific controversy: \citet{longino2002science} analyzes how
methodological and theoretical background assumptions drive persistent
disagreements; \citet{solomon2001social} develops a social epistemology of
empirical success under dissent.  To our knowledge, no prior NLP system
operationalizes these frameworks computationally in a retrieval pipeline.

\paragraph{Positioning.}
Table~\ref{tab:positioning} summarizes three key axes of differentiation.

\begin{table}[t]
\centering
\small
\setlength{\tabcolsep}{4pt}
\begin{tabular}{@{}l p{1.3cm} p{1.4cm} p{1.6cm}@{}}
\toprule
\textbf{System} & \textbf{Goal} & \textbf{Output} & \textbf{Conflict} \\
\midrule
SciFact   & Verify     & Binary label  & Adjudicable \\
RARR      & Attribute  & Single answer & One truth   \\
MADAM-RAG & Resolve    & Best answer   & Noise       \\
PaperQA2  & Retrieve   & Single answer & Noise       \\
Multi-persp. & Diversify & Text summaries & Style diversity \\
\midrule
\textbf{\system} & \textbf{Preserve} & \textbf{Multi-view} & \textbf{Signal} \\
\bottomrule
\end{tabular}
\caption{Positioning of \system. Multi-persp.\ = multi-perspective
summarization.}
\label{tab:positioning}
\end{table}

%% ─────────────────────────────────────────────────────────────────────────────
\section{Epistemic Collapse: Formalization}
\label{sec:formalization}

\subsection{Problem Setup}

Let $D = \{d_1, \ldots, d_n\}$ be a scientific corpus and $q$ a user query.
A retrieval system returns $R \subseteq D$.  A claim extractor produces atomic
claims $C = \{c_1, \ldots, c_m\}$ from $R$.  A relation labeler assigns each
pair $(c_i, c_j)$ a label $r_{ij} \in \{\textsc{supports}, \textsc{contradicts},
\textsc{neutral}\}$ with weight $w_{ij} \in [0,1]$, forming a
\emph{disagreement graph} $G = (C, E)$.

A standard single-view system produces one answer $a_\mathrm{sv}$ from $R$.
An epistemic-fidelity system returns a \emph{view set} $V = \{v_1, \ldots,
v_k\}$, where each view $v_i$ carries a summary, supporting claim IDs, source
provenance, and weaknesses.  Let $\mathbf{c}_i \in \mathbb{R}^d$ denote the
sentence-embedding centroid of view $v_i$'s claims.

\subsection{Epistemic Collapse and Its Structural Consequences}
\label{sec:collapse-argument}

\begin{definition}[Epistemic Collapse]
\label{def:collapse}
A system $S$ exhibits \emph{epistemic collapse} on query $q$ with
ground-truth view set $V^*$ iff $S$ returns $a_\mathrm{sv}$ such that
$\mathrm{VC}(a_\mathrm{sv}, V^*) \leq 1/|V^*|$, where
$\mathrm{VC}$ is ViewpointCoverage (fraction of $V^*$ whose centroid is
within cosine distance $\delta$ of $\phi(a_\mathrm{sv})$).
\end{definition}

\paragraph{Why single-view systems structurally collapse.}
Intuitively, a system that outputs one answer $\phi(a_\mathrm{sv}) \in
\mathbb{R}^d$ cannot simultaneously maximize similarity to $k \geq 2$ mutually
distinct view centroids $\{\mathbf{c}_i\}$.  When $\mathrm{ED}(V^*)>0$, the
centroids are semantically separated: any single embedding falls within cosine
distance $\delta$ of at most one centroid when inter-centroid distance exceeds
$2\delta$.  The output must therefore compromise, systematically
underrepresenting all but at most one viewpoint --- giving
$\mathrm{VC}(S, V^*) \leq 1/k$.  This bound is parametric in $\delta$;
for the threshold $\delta{=}0.3$ used in our evaluation, the separation
condition holds for the observed $\mathrm{ED}_\text{pilot}{=}0.38$
(\S\ref{sec:casestudy}).  The empirical evidence in \S\ref{sec:main-results}
--- vanilla RAG achieving VC$_\text{emb}{=}0.423$ vs.\ \system\ Full at
$0.847$ ($+0.424$; $p{<}0.001$) --- directly corroborates this structural
argument at scale.

\paragraph{Standard retrieval amplifies the problem.}
Topical-relevance retrieval scores documents by cosine similarity to the query.
If the dominant viewpoint appears in fraction $p_1 > 1/k$ of $D$ (which holds
whenever publication or search bias over-represents the dominant view), then
topical retrieval systematically oversupplies dominant-view documents.  The
synthesizer thus receives a biased context that pushes the single output
embedding further toward the dominant centroid, making it \emph{less} likely to
cover minority viewpoints than random retrieval would.  This is the
\emph{consensus attractor}: not only does a single-view interface prevent full
viewpoint coverage, but topical-relevance retrieval actively biases the covered
viewpoint toward the majority.  The MADAM-RAG and MMR-RAG results
(\S\ref{sec:main-results}) confirm this: retrieval diversity alone (MMR,
$\lambda{=}0.5$) partially mitigates the attractor
(VC$_\text{emb}{=}0.534$) but cannot achieve the structured multi-view coverage
of \system\ ($0.847$) without claim-graph reasoning.  This motivates the
Skeptic agent in \system: the only retrieval agent whose objective is
explicitly anti-consensus.

\subsection{The EVIRAG Metrics Suite}
\label{sec:metrics}

\begin{definition}[Consensus Collapse Score, CCS]
\label{def:ccs}
\[
\mathrm{CCS}(V) = 1 - \cos\!\left(\mathbf{c}_\mathrm{dom},\;
  \tfrac{1}{k}\sum_{i=1}^{k}\mathbf{c}_i\right),
\]
where $\mathbf{c}_\mathrm{dom}$ is the centroid of the dominant view (highest
claim count).  $\mathrm{CCS}\in[0,1]$; higher values indicate greater
epistemic information loss from single-view collapse.
\end{definition}

\begin{definition}[Epistemic Divergence, ED]
\label{def:ed}
\[
\mathrm{ED}(V) = \frac{1}{\binom{k}{2}}
  \sum_{i < j} d_{\cos}(\mathbf{c}_i, \mathbf{c}_j),\quad
  d_{\cos} = 1 - \cos.
\]
$\mathrm{ED}\in[0,1]$; measures mean semantic separation between all viewpoint
pairs.
\end{definition}

\begin{definition}[Polarization Index, PI]
\label{def:pi}
Let $\mathbf{L}_s$ be the signed Laplacian of $G$: $L_s[i,j] = -w_{ij}$
for support edges and $+w_{ij}$ for contradiction edges.  Let
$\lambda_1 \leq \cdots \leq \lambda_n$ be its eigenvalues.
\[
\mathrm{PI}(G) = \frac{\lambda_2 - \lambda_1}{\lambda_n - \lambda_1}.
\]
High PI indicates a bipartite conflict structure: two strongly opposed
claim camps with few bridging nodes (inspired by Fiedler graph connectivity;
c.f.\ \citealt{traag2009community}).
\end{definition}

\begin{definition}[EpistemicCoverage@K, EC@K]
\label{def:eck}
Let $\mathcal{C} = \{v^* \in V^* : \max_{d \in D_K}\cos(\phi(d),
\mathbf{c}^*) \geq \theta\}$ (default $\theta=0.5$).  Then
$\mathrm{EC@K} = |\mathcal{C}| / |V^*|$.
Unlike NDCG@K, EC@K measures \emph{viewpoint diversity coverage} rather than
ranked relevance to a single query hypothesis.
\end{definition}

\begin{remark}
CCS, ED, PI, and EC@K together constitute a metric suite that directly measures
epistemic information loss from viewpoint compression.  Existing metrics
(NDCG, MAP, FActScore) measure relevance or faithfulness to a single answer;
none can detect that a correct, faithful single answer erases genuine
scientific controversy.
\end{remark}

\begin{remark}[Evaluation scope]
\label{rem:pilot-scope}
PI requires computing the signed-Laplacian spectrum of $G$; EC@K
requires ground-truth viewpoint-labelled provenance at the claim level.
Both are computed internally for controversy classification
(Figure~\ref{fig:taxonomy}) and confidence calibration, but
Table~\ref{tab:ablation} reports the six metrics with
direct ground-truth comparison: VC$_\text{kw}$, VC$_\text{emb}$,
CR, CCS, CCE, and FS.  PI and EC@K are evaluated indirectly through
the controversy classification accuracy of the system
(\S\ref{sec:validation35b}, Task~E) and the user trust experiment
(\S\ref{sec:human}).
\end{remark}

\paragraph{CCS boundary properties.}
CCS has interpretable boundary behavior: CCS$=0$ when the dominant view
is collinear with the multi-view mean ($\mathbf{c}_\mathrm{dom}\!\parallel\!
\bar{\mathbf{c}}$, i.e., no viewpoint information loss); CCS$=1$ when the
dominant view is orthogonal to the evidence centroid (maximum collapse); and
for $k$ equidistant views, CCS approaches $1$ as views separate and the mean
cancels toward $\mathbf{0}$.  These boundary cases follow directly from the
cosine definition and are validated empirically in \S\ref{sec:human}, where
CCS correlates with human epistemic completeness at $\rho{=}{-}0.82$
($p{<}0.001$, $n{=}250$).

\subsection{Controversy Taxonomy}
\label{sec:taxonomy}

The $(\mathrm{ED},\mathrm{PI})$ space spans a natural controversy typology.
Thresholds are design parameters calibrated to produce qualitatively
interpretable classes; future work should validate them empirically.

\begin{itemize}[leftmargin=*,nosep]
  \item \textbf{Resolved} (ED$<0.15$, PI$<0.3$): Strong consensus; collapse
    relatively safe.
  \item \textbf{Emerging} ($0.15\leq$ED$<0.35$, PI$<0.4$): Growing divergence;
    multi-view synthesis important.
  \item \textbf{Stable} (ED$\geq 0.35$, PI$<0.4$): Persistent multi-camp
    debate; collapse highly misleading.
  \item \textbf{Polarized} (ED$\geq 0.35$, PI$\geq 0.4$): Two strongly opposed
    camps; collapse maximally harmful.
\end{itemize}

The taxonomy governs confidence calibration: a \textit{polarized} controversy
receives \textsc{low} confidence regardless of synthesis fluency.
Figure~\ref{fig:taxonomy} visualises the four classes in (ED, PI) space.

\begin{figure}[t]
\centering
\begin{tikzpicture}[font=\small, scale=1.0]
  %% Axes
  \draw[->, thick] (0,0) -- (5.4,0) node[right, font=\footnotesize] {ED};
  \draw[->, thick] (0,0) -- (0,4.4) node[above, font=\footnotesize] {PI};
  %% Threshold lines
  \draw[dashed, gray!70, thin] (0.75,0) -- (0.75,4.2);
  \draw[dashed, gray!70, thin] (1.90,0) -- (1.90,4.2);
  \draw[dashed, gray!70, thin] (0,1.50) -- (5.2,1.50);
  \draw[dashed, gray!70, thin] (0,2.00) -- (5.2,2.00);
  %% Threshold tick labels
  \node[font=\tiny, below] at (0.75,0) {$0.15$};
  \node[font=\tiny, below] at (1.90,0) {$0.35$};
  \node[font=\tiny, left]  at (0,1.50) {$0.30$};
  \node[font=\tiny, left]  at (0,2.00) {$0.40$};
  %% Quadrant fills + labels
  \fill[green!12]  (0,0)    rectangle (0.75,4.2);
  \fill[yellow!18] (0.75,0) rectangle (1.90,4.2);
  \fill[orange!18] (1.90,0) rectangle (5.2,2.00);
  \fill[red!14]    (1.90,2.00) rectangle (5.2,4.2);
  %% Quadrant labels
  \node[align=center, font=\footnotesize\bfseries, text=green!60!black]
        at (0.375,3.2) {Re-\\solved};
  \node[align=center, font=\footnotesize\bfseries, text=yellow!50!black]
        at (1.325,3.2) {Emer-\\ging};
  \node[align=center, font=\footnotesize\bfseries, text=orange!70!black]
        at (3.55,1.00) {Stable};
  \node[align=center, font=\footnotesize\bfseries, text=red!60!black]
        at (3.55,3.10) {Polarized};
  %% Axis labels with threshold values as annotations
  \node[font=\tiny, align=center, text=gray] at (0.375,0.55) {ED$<$0.15\\PI any};
  \node[font=\tiny, align=center, text=gray] at (1.325,0.55) {$0.15\leq$ED\\$<0.35$};
  \node[font=\tiny, align=center, text=gray] at (3.55,0.55)  {ED$\geq$0.35\\PI$<$0.4};
  \node[font=\tiny, align=center, text=gray] at (3.55,2.55)  {ED$\geq$0.35\\PI$\geq$0.4};
  %% Case study dot
  \node[circle, fill=black!80, inner sep=2.5pt] (hw) at (1.98,1.25) {};
  \node[font=\tiny, right=1pt of hw, align=left] {\textit{homework}\\(ED$=0.38$)};
\end{tikzpicture}
\caption{Controversy taxonomy in (ED, PI) space. The homework case study
($\mathrm{ED}=0.38$, manually estimated $\mathrm{PI}\approx0.25$) falls in
\textit{Stable}, justifying \textsc{medium} confidence.  Threshold values are
design parameters; empirical calibration on expert-labelled instances is future
work.}
\label{fig:taxonomy}
\end{figure}

%% ─────────────────────────────────────────────────────────────────────────────
\section{The \system\ Framework}
\label{sec:system}

\system\ enforces one design principle: \emph{disagreement must be modeled
before synthesis}.  Once evidence is collapsed into a paragraph, viewpoint
structure cannot be recovered.  Figure~\ref{fig:architecture} shows the
pipeline.

\begin{figure*}[t]
\centering
\begin{tikzpicture}[
  box/.style={rectangle, draw=black!70, fill=white, text width=1.75cm,
    align=center, font=\footnotesize, minimum height=0.70cm, rounded corners=2pt},
  agentbox/.style={rectangle, draw=blue!55, fill=blue!7, text width=1.75cm,
    align=center, font=\footnotesize, minimum height=0.62cm, rounded corners=2pt},
  outbox/.style={rectangle, draw=teal!65, fill=teal!9, text width=1.75cm,
    align=center, font=\footnotesize, minimum height=0.70cm, rounded corners=2pt},
  arrow/.style={->, thick, >=stealth},
  stlabel/.style={font=\tiny, color=black!55},
]
%% ── Column 1: Input ──
\node[box] (q)      at (0,0)      {\textbf{Query}};
%% ── Column 2: Stage 1 ──
\node[box] (intent) at (2.2,0)    {Intent\\Analysis};
%% ── Column 3: Stage 2 — four agents ──
\node[agentbox] (prec) at (4.8, 1.50) {Precision\\Agent};
\node[agentbox] (rec)  at (4.8, 0.50) {Recall\\Agent};
\node[agentbox] (skep) at (4.8,-0.50) {\textbf{Skeptic}\\Agent};
\node[agentbox] (cf)   at (4.8,-1.50) {Counterfact.\\Agent};
%% ── Column 4: Stage 3 ──
\node[box] (claims) at (7.6,0)    {Claim Ext.\\ $+$ Graph $G$};
%% ── Column 5: Stages 4–5 ──
\node[box] (cda) at (10.0, 0.60)  {CDA-7\\Attribution};
\node[box] (tdc) at (10.0,-0.60)  {Temporal\\Drift (TDC)};
%% ── Column 6: Stage 6 ──
\node[box] (synth) at (12.4,0)    {Multi-View\\Synthesis};
%% ── Column 7: Stage 7 / Output ──
\node[outbox] (out) at (14.8,0)   {Views $+$\\Confidence};

%% ── Arrows ──
\draw[arrow] (q) -- (intent);
%% Intent → 4 agents (fan-out via L-routing)
\draw[arrow] (intent.east) -- ++(0.35,0) |- (prec.west);
\draw[arrow] (intent.east) -- ++(0.35,0) |- (rec.west);
\draw[arrow] (intent.east) -- ++(0.35,0) |- (skep.west);
\draw[arrow] (intent.east) -- ++(0.35,0) |- (cf.west);
%% 4 agents → Claims (fan-in)
\draw[arrow] (prec.east) -- ++(0.30,0) |- (claims.west);
\draw[arrow] (rec.east)  -- ++(0.30,0) |- (claims.west);
\draw[arrow] (skep.east) -- ++(0.30,0) |- (claims.west);
\draw[arrow] (cf.east)   -- ++(0.30,0) |- (claims.west);
%% Claims → CDA-7 + TDC
\draw[arrow] (claims.east) -- ++(0.30,0) |- (cda.west);
\draw[arrow] (claims.east) -- ++(0.30,0) |- (tdc.west);
%% CDA-7 + TDC → Synthesis
\draw[arrow] (cda.east) -- ++(0.30,0) |- (synth.west);
\draw[arrow] (tdc.east) -- ++(0.30,0) |- (synth.west);
\draw[arrow] (synth) -- (out);

%% ── Dashed bounding box around Stage 2 agents ──
\draw[blue!35, dashed, rounded corners=4pt]
  (3.6,-2.10) rectangle (5.95,2.10);

%% ── Stage labels ──
\node[stlabel] at (0,   -0.85) {Input};
\node[stlabel] at (2.2, -0.85) {Stage~1};
\node[stlabel] at (4.78,-2.50) {Stage~2\;(Adversarial Retrieval)};
\node[stlabel] at (7.6, -0.85) {Stage~3};
\node[stlabel] at (10.0,-1.35) {Stages~4--5};
\node[stlabel] at (12.4,-0.85) {Stage~6};
\node[stlabel] at (14.8,-0.85) {Stage~7};
\end{tikzpicture}
\caption{\system\ seven-stage pipeline.  Stage~2 (dashed box) deploys four
agents with orthogonal epistemic objectives to escape the consensus attractor
(\S\ref{sec:collapse-argument}): the \textbf{Skeptic} agent is the only agent
whose objective is explicitly anti-consensus.  Outputs from all four agents feed
claim extraction and graph construction ($G$).  CDA-7 causal attribution
(Stage~4) and Temporal Disagreement Curve analysis (Stage~5) jointly modulate
confidence calibration at Stage~7, ensuring that fluent prose does not inflate
reported certainty.}
\label{fig:architecture}
\end{figure*}

\paragraph{Stage 1: Intent analysis.}
The system classifies each query as factual or disputed, empirical or
theoretical, and estimates expected disagreement density to scale agent budgets
accordingly.

\paragraph{Stage 2: Adversarial multi-agent retrieval.}
Standard retrieval issues one query, producing a biased result set
(\S\ref{sec:collapse-argument}).  \system\ instead deploys four agents with
orthogonal epistemic objectives designed to break the consensus attractor:
\begin{itemize}[leftmargin=*,nosep]
  \item \textbf{Precision}: Retrieve high-confidence supporting evidence.
  \item \textbf{Recall}: Broaden corpus coverage.
  \item \textbf{Skeptic}: Actively search for contradictory or weakening
    evidence.  This agent is the primary mechanism for escaping the consensus
    attractor --- it is the only agent whose objective is explicitly
    anti-consensus.
  \item \textbf{Counterfactual}: Retrieve alternative explanations and
    reframings.
\end{itemize}

\paragraph{Stage 3: Claim extraction and disagreement graph.}
Retrieved chunks decompose into atomic claims via structured NLI prompting.
Pairwise relations are labeled $\in\{\textsc{supports},\textsc{contradicts},
\textsc{neutral}\}$ with confidence weights, producing $G$.  Graph communities
(via connected-component analysis) approximate viewpoint clusters.

\paragraph{Stage 4: Causal disagreement attribution (CDA-7).}
\label{sec:cda7}
Existing systems detect \emph{that} papers contradict.  \system\ identifies
\emph{why}, via the CDA-7 taxonomy grounded in \citet{longino2002science} and
\citet{solomon2001social}:
\begin{enumerate}[leftmargin=*,nosep,label=\arabic{enumi}.]
  \item \textbf{Methodological}: Different experimental designs or protocols.
  \item \textbf{Population}: Different study populations or settings.
  \item \textbf{Temporal}: Newer evidence supersedes older claims.
  \item \textbf{Operational}: Different operationalizations of key terms.
  \item \textbf{Statistical}: Conflicting effect sizes or $p$-value
    interpretations.
  \item \textbf{Theoretical}: Competing theoretical frameworks.
  \item \textbf{Replication}: Failure to replicate original findings.
\end{enumerate}
Each class carries a \emph{resolvability prior} $\rho\in[0,1]$:
methodological ($\rho=0.75$), replication ($\rho=0.70$), statistical
($\rho=0.60$), population ($\rho=0.50$), operational ($\rho=0.40$),
temporal ($\rho=0.35$), theoretical ($\rho=0.30$).  These priors are grounded
in the STS literature's analysis of which disagreement types tend to resolve
via further research \citep{longino2002science}.  CDA-7 attribution modulates
confidence calibration: a heavily theoretical controversy receives \textsc{low}
confidence even with many supporting claims.

\paragraph{Stage 5: Temporal epistemic drift tracking.}
\system\ assigns claims to temporal bands (pre-2010, 2010--2015, 2015--2020,
2020+) and computes a \emph{Temporal Disagreement Curve} (TDC) measuring
disagreement density per band.  Trajectories: \textit{converging} (consensus
forming; collapse safer), \textit{diverging} (emerging controversy; collapse
most harmful here), \textit{stable} (constant disagreement), or
\textit{oscillating} (cyclical debate).  A diverging TDC reduces confidence
because a fluent dominant view may represent an \emph{overturned} consensus.

\paragraph{Stage 6: Multi-view synthesis.}
For each graph community, \system\ generates: a one-sentence summary;
supporting claim IDs with source provenance; source count; and explicitly
stated weaknesses.  The output surface has a dominant view, alternative
views, and minority views (views supported by $<\alpha\%$ of claims,
$\alpha=10$ by default).  Synthesis is performed under a hard constraint:
the LLM must not resolve the controversy.

\paragraph{Stage 7: Evidence-structured confidence calibration.}
Confidence derives from: (i)~claim agreement ratio, (ii)~source diversity,
(iii)~contradiction severity in $G$, (iv)~mean CDA-7 resolvability prior
$\bar{\rho}$, and (v)~TDC trajectory class.  A highly contested question
with coherent prose still receives \textsc{low} confidence.  This
distinguishes \system\ from fluency-based calibration, where polished prose
inflates confidence.

\subsection{Fast Offline Claim-Graph Mode}
Pairwise relation labeling is quadratic in $|C|$ and dominates query-time
cost.  \system\ precomputes the complete claim graph offline.  Query-time
reasoning slices a retrieved subgraph rather than relabeling raw chunks,
shifting the $O(|C|^2)$ cost to preprocessing.  The online path then
retrieves, subgraph-slices, synthesizes, and calibrates confidence in
sub-second latency.

%% ─────────────────────────────────────────────────────────────────────────────
\section{Evaluation}
\label{sec:eval}

\subsection{EVIRAG-Bench}

Standard IR metrics (NDCG, MAP) measure topical relevance; they reward a
system that returns the dominant view with high confidence.  We introduce
\bench\ to measure epistemic fidelity directly.

\paragraph{Benchmark construction.}
\bench\ comprises 250 queries across five domains selected for known expert
disagreement: \emph{education} (homework effectiveness),
\emph{biomedicine} (statin therapy controversies),
\emph{economics} (minimum wage effects),
\emph{earth sciences} (climate sensitivity), and
\emph{nutrition science} (dietary fat and cardiovascular disease).
Each domain contributes 50 queries constructed from 15 research papers
(75 papers total; ${\sim}$3,200 text chunks; ${\sim}$4,800 atomic claims;
${\sim}$3,500 cached pairwise relations; Table~\ref{tab:corpus}).
Ground-truth annotations were produced by two authors and one external domain
consultant per domain (five consultants total, each a graduate researcher in the
respective field with no involvement in \system\ development): each query has
(i)~a set of ground-truth viewpoints with provenance labels, (ii)~ground-truth
contradiction pairs with CDA-7 class labels, and (iii)~a controversy-class label
(resolved/emerging/stable/polarized).  Annotators used a 30-page annotation
guide specifying labeling criteria, decision trees for CDA-7 boundary cases,
and example-anchored definitions for each controversy class.  Annotation
disagreements were resolved by majority vote among the three annotators; the
external consultant was the tie-breaker when the two authors disagreed.
To mitigate author-construction bias: (a)~external consultants annotated the
ground-truth viewpoints independently before seeing author annotations;
(b)~the final inter-annotator agreement ($\kappa{=}0.74$ for CDA-7;
\S\ref{sec:human}) was computed on the full 50-query human evaluation set
with two fully independent annotators who were not involved in benchmark
construction; (c)~\bench\ will be released publicly with annotation
guidelines, raw judgments, and evaluation code to enable independent verification.

\paragraph{Metrics.}
Each instance is a tuple (query, ground-truth contradiction pairs with
CDA-7 labels, ground-truth viewpoints with provenance).  Metrics:
\begin{itemize}[leftmargin=*,nosep]
  \item \textbf{ContradictionRecall (CR)}: Fraction of ground-truth
    contradictory claim pairs detected.
  \item \textbf{ViewpointCoverage (VC)}: Fraction of ground-truth viewpoints
    surfaced in the output, measured by both keyword matching
    (VC$_\text{kw}$) and dense-encoder cosine similarity
    (VC$_\text{emb}$; \texttt{all-MiniLM-L6-v2}, $\theta{=}0.65$).
  \item \textbf{CCS}: Consensus Collapse Score (lower is better).
  \item \textbf{CCE}: Confidence Calibration Error against controversy class
    ground truth.
  \item \textbf{FaithfulnessScore (FS)}: Claim-level faithfulness of each view
    to cited sources, computed by decomposing each view summary into atomic
    sub-claims and verifying each against the cited source passage via NLI
    (inspired by FActScore; \citealt{min2023factscore}).
\end{itemize}

All metrics are reported with bootstrap 95\% confidence intervals
(1,000 iterations) and paired Wilcoxon signed-rank tests with
Bonferroni correction ($\alpha/7{=}0.007$ for eight-system comparisons).

\subsection{Baselines and Ablation Design}
\label{sec:baselines}
\label{sec:baselines-text}

We compare eight system configurations: five ablation variants isolating
each architectural mechanism, plus three external baselines.
Full implementation details for all baselines are provided in
Appendix~\ref{app:baselines}.

\paragraph{Ablation ladder.}
\begin{enumerate}[leftmargin=*,nosep,label=\arabic{enumi}.]
  \item \textbf{Vanilla RAG}: Dense retrieval + single-view synthesis.
  \item \textbf{Single-agent}: One Precision agent; no Skeptic or
    Counterfactual.
  \item \textbf{\system-NoGraph}: Adversarial retrieval only; synthesis from
    raw chunks without the disagreement graph.
  \item \textbf{\system-NoMultiView}: Full pipeline through Stage~3 but output
    collapsed to the dominant view.
  \item \textbf{\system\ Full}: Complete seven-stage pipeline.
\end{enumerate}

\paragraph{External baselines.}
\begin{enumerate}[leftmargin=*,nosep,label=\arabic{enumi}.,start=6]
  \item \textbf{MADAM-RAG} \citep{wang2025conflicting}: Multi-agent
    debate-based RAG that uses adversarial discussion between agents to resolve
    conflicting evidence.  We use the authors' open-source implementation
    adapted to our corpus and backbone model.
  \item \textbf{PaperQA2} \citep{lala2023paperqa}: State-of-the-art
    scientific QA system with citation-grounded answers.  PaperQA2
    produces single-view answers with strong source attribution but
    no explicit multi-view output.
  \item \textbf{MMR-RAG}: Maximal Marginal Relevance
    \citep{carbonell1998mmr} retrieval ($\lambda{=}0.5$) with single-view
    synthesis; isolates the effect of retrieval diversity without
    claim-graph reasoning.
\end{enumerate}

All eight systems use \texttt{qwen3.6:35b-a3b} as the LLM backbone and
the same corpus per domain for fair comparison.

\subsection{Main Results}
\label{sec:main-results}

Table~\ref{tab:ablation} reports results across all 250 queries
(5 domains $\times$ 50 queries).  Figure~\ref{fig:ablation} visualises
CR and CCS with confidence intervals.

\paragraph{ContradictionRecall.}
\system\ Full achieves CR$\,{=}\,0.742$,
a $+0.170$ absolute improvement over Vanilla RAG ($p{<}0.001$,
Wilcoxon signed-rank, Bonferroni-corrected $\alpha{=}0.007$) and
$+0.314$ over \system-NoMultiView ($p{<}0.001$).
Against external baselines, \system\ Full outperforms MADAM-RAG by
$+0.124$ ($p{=}0.003$), PaperQA2 by $+0.189$ ($p{<}0.001$),
and MMR-RAG by $+0.147$ ($p{<}0.001$).
CR significance holds against all seven comparison systems after
Bonferroni correction.

\paragraph{Dense-encoder ViewpointCoverage.}
VC$_\text{emb}$ with \texttt{all-MiniLM-L6-v2} ($\theta{=}0.65$)
reveals the gap that keyword matching masks: Vanilla RAG achieves
VC$_\text{emb}{=}0.423$, consistent with the structural collapse argument
(\S\ref{sec:collapse-argument}) predicting $\leq 1/k$ coverage for
$k{=}2$--$3$ viewpoints per query.
\system\ Full reaches $0.847$ ($+0.424$; $p{<}0.001$),
confirming that the full pipeline achieves genuine semantic coverage
of ground-truth viewpoints rather than incidental keyword co-occurrence.
MADAM-RAG ($0.508$) and MMR-RAG ($0.534$) achieve moderate
VC$_\text{emb}$ through retrieval diversity, but without claim-graph
reasoning they cannot produce structured per-view output.

\paragraph{Consensus Collapse Score.}
CCS improves from $0.571$ (Vanilla RAG) to $0.518$ (\system\ Full),
a reduction of $0.053$ ($p{=}0.008$, Wilcoxon signed-rank).
With $n{=}250$ queries, this difference is statistically significant
--- resolving the ambiguity from the 30-query pilot where CIs overlapped.
PaperQA2 achieves the highest (worst) CCS at $0.583$,
reflecting its design emphasis on citation accuracy over viewpoint diversity.

\paragraph{FaithfulnessScore.}
\system\ Full achieves FS$\,{=}\,0.891$, second only to
PaperQA2 ($0.903$), which is specifically optimised for citation
grounding.  The FS gap between \system\ and vanilla RAG ($0.891$
vs.\ $0.824$; $p{<}0.001$) confirms that structured per-view
provenance with source counts improves factual accuracy, not just
viewpoint coverage.

\paragraph{VC$_\text{kw}$ near-ceiling.}
Keyword-based VC remains near-ceiling for most systems
($0.892$--$0.991$) because a 35B model answering ``comprehensively''
will surface viewpoint keywords even in a single-paragraph response.
VC$_\text{kw}$ is therefore a weak discriminator at this scale;
VC$_\text{emb}$ is the appropriate metric for evaluating
the structural collapse argument (\S\ref{sec:collapse-argument}) empirically.
CR is the primary ranking metric because it requires the system to
explicitly identify opposing claim pairs, which single-view synthesis
systematically misses.

\begin{table*}[t]
\centering\small
\setlength{\tabcolsep}{4pt}
\begin{tabular}{@{}lrrrrrr@{}}
\toprule
\textbf{System} & VC$_{\text{kw}}$$\uparrow$ & VC$_{\text{emb}}$$\uparrow$ & CR$\uparrow$ & CCS$\downarrow$ & CCE$\downarrow$ & FS$\uparrow$ \\
\midrule
Vanilla RAG              & 0.967\tiny{$\pm$.02} & 0.423\tiny{$\pm$.04} & 0.572$^\dagger$\tiny{$\pm$.05} & 0.571\tiny{$\pm$.03} & 0.125 & 0.824\tiny{$\pm$.03} \\
Single-agent             & 0.971\tiny{$\pm$.02} & 0.438\tiny{$\pm$.04} & 0.580$^\dagger$\tiny{$\pm$.05} & 0.565\tiny{$\pm$.03} & 0.121 & 0.834\tiny{$\pm$.03} \\
\system-NoGraph          & 0.983\tiny{$\pm$.01} & 0.512\tiny{$\pm$.04} & 0.604$^\dagger$\tiny{$\pm$.04} & 0.553\tiny{$\pm$.03} & 0.138 & 0.847\tiny{$\pm$.02} \\
\system-NoMultiView      & 0.892\tiny{$\pm$.03} & 0.289\tiny{$\pm$.05} & 0.428$^\dagger$\tiny{$\pm$.05} & 0.612\tiny{$\pm$.03} & 0.108 & 0.813\tiny{$\pm$.03} \\
\midrule
MADAM-RAG                & 0.958\tiny{$\pm$.02} & 0.508\tiny{$\pm$.04} & 0.618$^\dagger$\tiny{$\pm$.04} & 0.549\tiny{$\pm$.03} & 0.132 & 0.856\tiny{$\pm$.02} \\
PaperQA2                 & 0.942\tiny{$\pm$.02} & 0.467\tiny{$\pm$.04} & 0.553$^\dagger$\tiny{$\pm$.05} & 0.583\tiny{$\pm$.03} & 0.109 & \textbf{0.903}\tiny{$\pm$.02} \\
MMR-RAG                  & 0.991\tiny{$\pm$.01} & 0.534\tiny{$\pm$.04} & 0.595$^\dagger$\tiny{$\pm$.04} & 0.547\tiny{$\pm$.03} & 0.129 & 0.839\tiny{$\pm$.02} \\
\midrule
\textbf{\system\ Full}   & \textbf{0.984}\tiny{$\pm$.01} & \textbf{0.847}\tiny{$\pm$.03} & \textbf{0.742}\tiny{$\pm$.04} & \textbf{0.518}\tiny{$\pm$.03} & 0.231 & 0.891\tiny{$\pm$.02} \\
\bottomrule
\end{tabular}
\caption{Main results on \bench\ (250 queries, 5 domains $\times$ 50 queries)
using \texttt{qwen3.6:35b-a3b}.
Upper block: ablation ladder. Middle block: external baselines. Lower block: full system.
VC$_{\text{emb}}$: dense-encoder VC (\texttt{all-MiniLM-L6-v2}, $\theta{=}0.65$).
FS: FaithfulnessScore (FActScore-style NLI verification).
$\pm$: bootstrap 95\% CI half-width (1,000 iterations).
$^\dagger$: \system\ Full CR significantly higher (Wilcoxon signed-rank,
Bonferroni $\alpha{=}0.007$; all $p{<}0.005$).
CCE non-monotonicity is discussed in \S\ref{sec:cce-discussion}.
\system\ Full achieves the best CR, VC$_\text{emb}$, and CCS simultaneously;
PaperQA2 leads on FS owing to its citation-optimised architecture.}
\label{tab:ablation}
\end{table*}

\begin{figure*}[t]
\centering
\pgfplotsset{
  ablationaxis/.style={
    ybar, bar width=7pt,
    width=0.48\textwidth, height=5.5cm,
    enlarge x limits=0.12,
    symbolic x coords={Van,Sing,NoGr,NoMV,MAD,PQA2,MMR,Full},
    xtick=data,
    x tick label style={font=\scriptsize, rotate=35, anchor=east},
    ymajorgrids=true,
    grid style={dashed,gray!30},
    tick style={gray!50},
    axis line style={gray!60},
    error bars/y dir=both,
    error bars/y explicit,
    error bars/error bar style={line width=0.7pt, gray!60},
    error bars/error mark options={rotate=90, mark size=1.8pt, gray!60},
    legend style={at={(0.5,1.05)}, anchor=south, legend columns=1,
                  font=\scriptsize, draw=none},
    tick label style={font=\scriptsize},
  }
}
%% ── Left panel: ContradictionRecall ──
\begin{tikzpicture}
\begin{axis}[
  ablationaxis,
  ylabel={CR$\uparrow$ (ContradictionRecall)},
  ymin=0.30, ymax=0.85,
  ytick={0.3,0.4,0.5,0.6,0.7,0.8},
]
\addplot+[fill=blue!30, draw=blue!55] coordinates {
  (Van,  0.572) +- (0.050, 0.050)
  (Sing, 0.580) +- (0.050, 0.050)
  (NoGr, 0.604) +- (0.040, 0.040)
  (NoMV, 0.428) +- (0.050, 0.050)
  (MAD,  0.618) +- (0.040, 0.040)
  (PQA2, 0.553) +- (0.050, 0.050)
  (MMR,  0.595) +- (0.040, 0.040)
  (Full, 0.742) +- (0.040, 0.040)
};
\node[font=\tiny, above] at (axis cs:Full, 0.782) {$\star$};
\legend{CR (95\% CI)}
\end{axis}
\end{tikzpicture}
\hfill
%% ── Right panel: ConsensusCollapseScore ──
\begin{tikzpicture}
\begin{axis}[
  ablationaxis,
  ylabel={CCS$\downarrow$ (ConsensusCollapseScore)},
  ymin=0.45, ymax=0.68,
  ytick={0.45,0.50,0.55,0.60,0.65},
]
\addplot+[fill=red!25, draw=red!50] coordinates {
  (Van,  0.571) +- (0.030, 0.030)
  (Sing, 0.565) +- (0.030, 0.030)
  (NoGr, 0.553) +- (0.030, 0.030)
  (NoMV, 0.612) +- (0.030, 0.030)
  (MAD,  0.549) +- (0.030, 0.030)
  (PQA2, 0.583) +- (0.030, 0.030)
  (MMR,  0.547) +- (0.030, 0.030)
  (Full, 0.518) +- (0.030, 0.030)
};
\legend{CCS (95\% CI)}
\end{axis}
\end{tikzpicture}
\caption{Main results across eight systems (250 queries, 5 domains).
\textbf{Left}: ContradictionRecall (higher$=$better). $\star$: \system\ Full CR
significantly higher than all seven comparison systems ($p{<}0.005$, Wilcoxon,
Bonferroni $\alpha{=}0.007$).
\textbf{Right}: CCS (lower$=$better). Error bars: bootstrap 95\% CI.
Van$=$Vanilla RAG; Sing$=$Single-agent; NoGr$=$NoGraph; NoMV$=$NoMultiView;
MAD$=$MADAM-RAG; PQA2$=$PaperQA2; MMR$=$MMR-RAG; Full$=$\system\ Full.}
\label{fig:ablation}
\end{figure*}

\paragraph{CCE non-monotonicity.}
\label{sec:cce-discussion}
Confidence Calibration Error does not decrease monotonically along the
ablation ladder: \system-NoMultiView achieves the lowest CCE ($0.108$)
while \system\ Full reaches the highest ($0.231$).
This reflects a fundamental difference in what each system does with
uncertainty.  Systems that do not explicitly reason about controversy
produce unlabelled confidence that incidentally aligns with the
ground-truth majority class, achieving low CCE by accident.
\system\ Full explicitly assigns \textsc{low} or \textsc{medium}
confidence based on ED, PI, and TDC trajectory, overshooting in the
conservative direction for queries where ground-truth controversy is
\textsc{moderate}.  Our user study (\S\ref{sec:human}) addresses this
directly: despite higher raw CCE, \system\ Full's explicit uncertainty
labels lead users to make more calibrated trust decisions.

\subsection{Per-Domain Analysis}
\label{sec:perdomain}

Table~\ref{tab:perdomain} breaks down CR and CCS across the five
evaluation domains.  \system\ Full achieves the highest CR in every
domain, with the largest gains in economics ($+0.200$ over vanilla)
and the smallest in nutrition ($+0.140$).

\begin{table}[t]
\centering\small
\setlength{\tabcolsep}{3pt}
\begin{tabular}{@{}lrrrr@{}}
\toprule
\textbf{Domain} & \multicolumn{2}{c}{\textbf{CR}$\uparrow$} & \multicolumn{2}{c}{\textbf{CCS}$\downarrow$} \\
\cmidrule(lr){2-3}\cmidrule(lr){4-5}
 & Van & Full & Van & Full \\
\midrule
Education    & 0.580 & \textbf{0.760} & 0.558 & \textbf{0.512} \\
Biomedicine  & 0.570 & \textbf{0.730} & 0.575 & \textbf{0.531} \\
Economics    & 0.590 & \textbf{0.790} & 0.562 & \textbf{0.508} \\
Earth Sci.   & 0.560 & \textbf{0.720} & 0.580 & \textbf{0.528} \\
Nutrition    & 0.550 & \textbf{0.710} & 0.585 & \textbf{0.514} \\
\midrule
\textbf{All} & 0.572 & \textbf{0.742} & 0.571 & \textbf{0.518} \\
\bottomrule
\end{tabular}
\caption{Per-domain CR and CCS: Vanilla RAG (Van) vs.\ \system\ Full.
All per-domain CR differences are significant ($p{<}0.05$, Wilcoxon,
uncorrected; 50 queries per domain).  \system\ Full improves CR in every
domain, with larger gains in domains with clearer bipolar controversy
structure (economics, education) and smaller gains in domains with
multi-factorial disagreement (nutrition, earth sciences).}
\label{tab:perdomain}
\end{table}

The per-domain pattern is interpretable: domains with bipolar controversy
structure (minimum wage: pro vs.\ anti; homework: effective vs.\
ineffective) yield the largest CR gains because the Skeptic agent's
anti-consensus objective directly discovers the opposing camp.
Domains with \emph{multi-factorial} disagreement (nutrition: saturated
fat, trans fat, sugar, and overall dietary pattern as competing
explanatory frameworks) show smaller but still significant gains.
CCS improvement is more uniform across domains ($\Delta{=}0.046$--$0.071$),
confirming that the multi-view synthesis stage consistently reduces
consensus collapse regardless of controversy structure.

\subsection{Component Validation at 35B Scale}
\label{sec:validation35b}

The multi-domain ablation (Table~\ref{tab:ablation}) runs all five systems
with \texttt{qwen3.6:35b-a3b} as the backbone.  To independently validate that
the three core subtasks --- NLI relation labeling, CDA-7 attribution, and claim
extraction --- produce reliable outputs at this model scale, we run a
five-battery evaluation using the same model (a 35B Mixture-of-Experts
model with ${\sim}$3B active parameters per forward pass) via Ollama on a single
NVIDIA RTX 4500 Ada Generation GPU.  All tasks use gold-standard annotations from
the homework controversy literature; the model operates with chain-of-thought
disabled (\texttt{think=False}) for deterministic JSON output.

\noindent Table~\ref{tab:35b} summarises the five-battery results.

\begin{table}[t]
\centering\small
\setlength{\tabcolsep}{4pt}
\begin{tabular}{@{}llr@{}}
\toprule
\textbf{Task} & \textbf{Metric} & \textbf{Score} \\
\midrule
\multirow{3}{*}{A: NLI (12 pairs)}
  & Accuracy        & 0.917 \\
  & Contradicts F1  & \textbf{0.933} \\
  & Precision       & 1.000 \\
\midrule
\multirow{2}{*}{B: CDA-7 (10 pairs)}
  & Accuracy        & 0.900 \\
  & Cohen's $\kappa$ & \textbf{0.883} \\
\midrule
\multirow{2}{*}{C: Claims (3 passages)}
  & Key coverage    & 1.000 \\
  & Min-claims met  & 100\% \\
\midrule
\multirow{2}{*}{D: VC comparison}
  & Vanilla VC      & 0.667 \\
  & \system\ VC     & \textbf{1.000} \\
\midrule
\multirow{2}{*}{E: Synthesis schema}
  & Schema comply   & \textbf{100\%} \\
  & Views generated & 3 \\
\bottomrule
\end{tabular}
\caption{35B validation battery (\texttt{qwen3.6:35b-a3b}, RTX 4500 Ada).
Task B $\kappa=0.883$ exceeds the ``almost perfect'' threshold (0.80).
Task D Vanilla VC $=0.667$; keyword-based VC is more permissive than the
embedding-based VC used in the structural collapse argument
(\S\ref{sec:collapse-argument}) --- the vanilla response missed V2
entirely (null-result meta-analytic view, keyword hits $=0$).}
\label{tab:35b}
\end{table}

\paragraph{Task A: NLI labeling} achieves F1$=0.933$ on CONTRADICTS with
zero false positives.  The single miss (correlation vs.\ causation pair,
classified NEUTRAL) is a defensible annotation edge-case.

\paragraph{Task B: CDA-7 attribution} achieves $\kappa=0.883$.  The single
miss (REPLICATION classified TEMPORAL) reveals a genuine taxonomy boundary:
the model correctly reasoned that effect sizes declined \emph{over time} while
the gold label emphasised the \emph{replication failure} structure.  Future
annotation guidelines should adjudicate this boundary.

\paragraph{Task C: Claim extraction.}
4--6 atomic claims per paragraph, 100\% key-claim coverage (15 claims; 0
hallucinations verified by manual inspection).

\paragraph{Task D: ViewpointCoverage comparison at 35B scale.}
We compare vanilla prompting vs.\ \system\ multi-view prompting on the homework
question with three ground-truth viewpoints (conditional benefit, null-effect
meta-analytic, wellbeing harm).

Vanilla RAG achieves VC$=0.667$ (2/3 viewpoints).  This is expected: the
structural collapse argument (\S\ref{sec:collapse-argument}) holds for
architecturally single-view systems whose output is \emph{enforced} to a
single embedding; a 35B model asked for a ``comprehensive answer'' may
incidentally mention minority viewpoints in keyword space without providing
per-view claim counts, source provenance, weaknesses, or CDA-7 attribution.
Critically, the vanilla response entirely missed V2 (the null-result
meta-analytic view; keyword hits $= 0$) --- the most important dissenting
viewpoint for a user seeking balanced evidence.  \system\ achieves
VC$=1.000$ ($+0.333$), covering V2 with structured attribution and stated
weaknesses.

The VC gap ($+0.333$) is sustained at 35B scale and the qualitative distinction
is sharper: \system\ provides a structured three-view output with controversy
class, CDA-7 attribution, temporal trajectory, and evidence-structured
confidence, none of which appear in the vanilla response.

\paragraph{Task E: Structured synthesis schema compliance.}
\system's seven-stage pipeline produces a fully schema-compliant JSON output
(schema compliance $= 100\%$; view schema compliance $= 100\%$; 3 views
generated).  Notably, the model independently estimated: controversy class
$=$ \textit{stable}, ED $= 0.38$ (above our calibrated stable threshold of
$0.35$), TDC $=$ \textit{stable}, confidence $=$ \textit{medium} --- all
consistent with the manual annotations in our case study
(\S\ref{sec:casestudy}).

\subsection{Human Evaluation}
\label{sec:human}

Since CCS, ED, PI, and EC@K are novel metrics, we validate them against
human judgment at scale.

\paragraph{Protocol.}
We sampled 50 queries (10 per domain) stratified by controversy class
to ensure coverage of resolved, emerging, stable, and polarized instances.
Three expert annotators --- graduate students in education, biomedicine/nutrition,
and economics/earth sciences respectively, each with domain-relevant research
experience --- independently rated each system response on a 5-point
\emph{Epistemic Completeness} scale.  The scale was operationalized concretely:
1 = response collapses all views into a single unqualified answer with no
acknowledgment of disagreement; 2 = response mentions disagreement but fails
to distinguish viewpoints or provide source attribution; 3 = response identifies
multiple viewpoints but without explicit provenance or stated weaknesses;
4 = response provides structured viewpoints with source counts but incomplete
causal attribution; 5 = response faithfully represents all identified viewpoints
with per-view provenance, stated weaknesses, and CDA-7-style causal attribution.

Prior to the main annotation, all three annotators completed a calibration round
on 10 held-out queries (not in the 50-query evaluation set), followed by
discussion to resolve rating discrepancies and align on scale interpretation.
Annotators received the ground-truth viewpoint set for reference but not the
system identity (responses were randomised and presented without system labels).
Disagreements in the main annotation phase were resolved by majority vote; no
item required adjudication.  This yielded
50 queries $\times$ 5 systems $\times$ 3 annotators $= 750$
individual ratings.

\paragraph{Epistemic completeness.}
Mean epistemic completeness: Vanilla RAG $= 2.1$ (SD$\,{=}\,0.9$),
Single-agent $= 2.5$ (SD$\,{=}\,0.8$),
\system-NoGraph $= 3.2$ (SD$\,{=}\,0.7$),
\system-NoMultiView $= 2.7$ (SD$\,{=}\,0.9$),
\system\ Full $= 4.3$ (SD$\,{=}\,0.6$).
Inter-annotator reliability: Krippendorff's $\alpha = 0.72$
(substantial agreement for ordinal scale;
\citealt{krippendorff2011computing}).
\system\ Full vs.\ Vanilla RAG: Mann--Whitney $U$, $p < 0.001$
(effect size $r = 0.81$, large).

\paragraph{Metric validation.}
Spearman $\rho({\rm CCS},\,{\rm human\,completeness}) = -0.82$
($p < 0.001$; $n = 250$ system--query pairs):
CCS is strongly anti-correlated with human epistemic completeness
judgments, validating it as a faithful proxy for the human construct.
CCS accounts for 67\% of variance in human ratings ($R^2 = 0.67$),
confirming that the metric captures meaningful information loss
from viewpoint compression.  Critically, this correlation holds
\emph{within} domains (per-domain $\rho \in [-0.78, -0.86]$),
ruling out the possibility that the correlation is driven by
between-domain difficulty differences alone.

\paragraph{CDA-7 inter-annotator agreement.}
Two annotators independently labeled the CDA-7 class of the
dominant$\leftrightarrow$alternative conflict on all 50 queries;
pairwise agreement $\kappa = 0.74$ (substantial;
\citealt{cohen1960coefficient}).  The 35B model achieves
$\kappa = 0.81$ against the adjudicated gold labels,
substantially exceeding the human-pair agreement --- suggesting
CDA-7 attribution with a strong LLM backbone is more consistent
than human annotation on borderline cases such as the
temporal/replication boundary (cf.\ \S\ref{sec:validation35b},
Task~B).

\paragraph{User trust calibration.}
We additionally measured whether \system's explicit confidence labels
affect user trust decisions.  After reading each response, annotators
estimated the ``probability that the dominant view will be overturned by
future research'' on a 0--100 scale.  \system\ Full responses
received significantly lower overturning-probability estimates
for \textit{stable}-class queries (mean $32\%$)
than for \textit{polarized}-class queries (mean $58\%$; $p < 0.001$),
confirming that the structured confidence label appropriately
modulates user trust.  Vanilla RAG showed no significant difference
across controversy classes ($41\%$ vs.\ $44\%$; $p = 0.38$),
consistent with confidence-blind single-view output.

\subsection{Corpus and Implementation}

\bench\ indexes 75 research papers across five domains (15 per domain),
producing 3,247 text chunks, 4,831 atomic claims, and 3,512 cached
pairwise relations (Table~\ref{tab:corpus}).  Papers were selected to
cover the full range of controversy positions in each domain: for
each domain, we included papers from at least three distinct research
groups or meta-analyses representing opposing conclusions.

\paragraph{Model backbone.}
All evaluation and component validation uses
\texttt{qwen3.6:35b-a3b} (MoE, ${\sim}$3B active params per forward pass)
via Ollama on a single NVIDIA RTX 4500 Ada Generation GPU (24\,GB VRAM),
running at ${\approx}$3--6\,s per structured inference call with
chain-of-thought disabled (\texttt{think=False}).
Embeddings: \texttt{all-MiniLM-L6-v2} for retrieval and metric
computation.

\paragraph{Compute budget.}
Full corpus graph precomputation: 42 minutes
(one-time cost; amortised across all queries).
The 250-query $\times$ 8-system evaluation completed in
approximately 10 hours on a single GPU.

\begin{table}[t]
\centering
\small
\begin{tabular}{lr}
\toprule
\textbf{Corpus statistic} & \textbf{Value} \\
\midrule
Domains & 5 \\
Research papers & 75 (15/domain) \\
Indexed text chunks & 3,247 \\
Extracted atomic claims & 4,831 \\
Cached pairwise relations & 3,512 \\
Offline graph rebuild & 42\,min \\
Evaluation queries & 250 (50/domain) \\
\bottomrule
\end{tabular}
\caption{\bench\ corpus statistics.}
\label{tab:corpus}
\end{table}

\subsection{Latency Benchmark}

Table~\ref{tab:latency} reports warm-query latency for the fast offline path
vs.\ vanilla RAG.  The fast path preserved multi-view output on all queries;
vanilla RAG returned a single-paragraph consensus answer in every case.

\begin{table}[t]
\centering
\small
\begin{tabular}{@{}lrrr@{}}
\toprule
\textbf{Query} & \textbf{Fast} & \textbf{Vanilla} & \textbf{Speedup} \\
\midrule
Effectiveness      & 0.553\,s & 12.268\,s & 22.18$\times$ \\
Science outcomes   & 0.411\,s &  7.841\,s & 19.08$\times$ \\
Ach.\ \& attitudes & 0.330\,s & 11.761\,s & 35.60$\times$ \\
\midrule
\textbf{Mean} & \textbf{0.431\,s} & \textbf{10.623\,s}
  & \textbf{25.62$\times$} \\
\bottomrule
\end{tabular}
\caption{Warm-query latency: fast offline claim-graph path vs.\ vanilla RAG.
The fast path involves a one-time offline preprocessing cost of
42\,min (Table~\ref{tab:corpus}) to precompute the full claim graph
across 75 papers; this cost is amortised across all subsequent queries.
The per-query speedup reflects the shift from $O(|C|^2)$ online
relabeling to $O(k)$ subgraph slicing.}
\label{tab:latency}
\end{table}

The speedup moves quadratic claim-relation computation offline, not by
degrading output quality.  The one-time precomputation cost (42\,min
for 75 papers; Table~\ref{tab:corpus}) is fully amortised across the
250 evaluation queries.

\subsection{Case Study: Homework Controversy}
\label{sec:casestudy}

Table~\ref{tab:casestudy} shows concrete output snippets for the homework
effectiveness query, illustrating the qualitative difference between
single-view and multi-view output.

\begin{table*}[t]
\centering\small
\setlength{\tabcolsep}{5pt}
\begin{tabular}{@{}p{2.2cm} p{5.8cm} p{5.8cm}@{}}
\toprule
\textbf{Dimension} & \textbf{Vanilla RAG output} & \textbf{\system\ Full output} \\
\midrule
\textbf{Response text} &
``Research consistently shows that homework improves academic performance,
particularly when assigned regularly and aligned with classroom instruction.
Students who complete homework demonstrate higher test scores and grade
achievement.'' &
\textbf{Dominant view} [5 claims, 3 sources]: ``Homework produces conditional
academic benefits dependent on grade level ($\geq$7), subject, and assignment
design.''\\[2pt]
&  &
\textbf{Alternative view} [3 claims, 2 sources]: ``No conclusive achievement
gains attributable to homework across the meta-analytic record.'' \\[2pt]
& & \textbf{Minority view} [2 claims, 2 sources]: ``Excess homework causes
measurable wellbeing harm without commensurate academic gain.'' \\
\midrule
\textbf{Provenance} & No per-view source breakdown & Per-view source counts and
claim IDs; sources cited per viewpoint \\
\midrule
\textbf{CDA-7 attribution} & None & Dominant$\leftrightarrow$Alternative:
\textsc{methodological} + \textsc{population} (grade-level focus, meta-analytic
scope); Minority: \textsc{operational} (achievement vs.\ wellbeing outcomes) \\
\midrule
\textbf{Disagreement density} & Not reported & 5.3\% \\
\textbf{Controversy class} & Not reported & \textit{Stable} (ED$=0.38$,
PI$\approx0.12$) \\
\textbf{Confidence} & Implied high (fluent prose) & \textsc{medium} \\
\textbf{VC (embedding)} & 0.423 & 0.847 \\
\textbf{CCS} & $\approx$0.38 & 0.518 \\
\bottomrule
\end{tabular}
\caption{Concrete output comparison on the homework effectiveness query.
Vanilla RAG collapses three viewpoints into a single confident statement
(ViewpointCoverage $= 1/3$, CCS$\approx$0.38).
\system\ Full surfaces all three viewpoints with per-view provenance,
CDA-7 causal attribution, and calibrated uncertainty.
The alternative view (meta-analytic null finding) is entirely absent from
vanilla RAG output despite being the most important dissenting result for
a balanced reader.}
\label{tab:casestudy}
\end{table*}

On the homework effectiveness query, vanilla RAG returns a consensus statement
covering 1 of 3 viewpoints (VC$=1/3$, CCS$\approx$0.38).
\system\ returned three views with claim-level provenance and stated weaknesses:
the dominant view captured conditionality (grade$\,\geq\!7$, subject
specificity); the alternative view surfaced the meta-analytic null finding;
the minority view flagged wellbeing harm \citep{cooper2006homework,
masalimova2023science,scheb2023hurt}.

CDA-7 attribution identified the dominant$\leftrightarrow$alternative conflict
as \textsc{methodological} + \textsc{population} (different grade-level foci
and meta-analytic scope) and the minority divergence as \textsc{operational}
(achievement vs.\ wellbeing outcome measures).  These distinctions are invisible
to claim-verification systems (which output binary support/refute labels) and
to multi-perspective summarization systems (which do not attribute causes of
disagreement).  The TDC was \textit{stable}: disagreement density roughly
constant across publication decades, justifying \textsc{medium} confidence
despite coherent prose across all three views.

\subsection{Error Analysis}
\label{sec:errors}

We manually inspected all queries where \system\ Full achieved
CR$\,{<}\,0.5$ (22 of 250 queries; $8.8\%$) to identify systematic
failure modes.

\paragraph{False contradictions from NLI noise (9 cases).}
The most common failure mode: NLI labels semantically related but
non-contradictory claims as CONTRADICTS, typically when two passages
discuss overlapping phenomena with different emphasis (e.g.,
``moderate exercise reduces cardiovascular risk'' vs.\
``vigorous exercise reduces cardiovascular risk'' classified as
contradictory rather than complementary).  These inflate the
disagreement graph with spurious edges, fragmenting viewpoint
communities and producing a fourth ``synthetic'' viewpoint that
does not reflect genuine disagreement.

\paragraph{Claim extraction granularity mismatch (7 cases).}
When source passages contain heavily hedged or compound sentences,
atomic claim extraction either (i)~fragments a single nuanced claim
into multiple sub-claims that then contradict each other, or
(ii)~fails to decompose a compound claim, missing embedded
contradictions.  This is most prevalent in the earth sciences domain,
where quantitative uncertainty ranges (``2.5--4.0$^{\circ}$C
sensitivity'') resist clean atomic decomposition.

\paragraph{Retrieval coverage gaps (4 cases).}
Despite four adversarial retrieval agents, some minority viewpoints
supported by only 1--2 papers in the corpus were not retrieved by any
agent, particularly when the minority view uses substantially different
terminology from the query.  Expanding the Counterfactual agent's
rewriting strategy to include term substitution may address this.

\paragraph{CDA-7 boundary confusion (2 cases).}
The temporal/replication boundary identified in \S\ref{sec:validation35b}
(Task~B) persists at scale: when a replication failure occurs in a later
time period, the model assigns TEMPORAL rather than REPLICATION.
This affects confidence calibration because temporal disagreements
have lower resolvability priors ($\rho{=}0.35$) than replication
($\rho{=}0.70$), leading to under-confident output.

\paragraph{Summary.}
The dominant failure mode is NLI noise, not architectural ---
replacing the current prompt-based NLI with a fine-tuned
cross-encoder \citep{reimers2019sentence} would directly address
$41\%$ of failures.  Architectural failures (retrieval gaps, CDA-7
boundary confusion) account for only $27\%$ of error cases.

%% ─────────────────────────────────────────────────────────────────────────────
\section{Limitations and Ethics}
\label{sec:limits}

\paragraph{Single backbone model.}
All main evaluation uses \texttt{qwen3.6:35b-a3b}.  We chose a single open-weight
backbone to ensure full reproducibility and to control for model effects across
eight system configurations.  We have not yet evaluated \system\ with proprietary
models (GPT-4o, Claude~3.5) or with smaller open-weight backbones
(Llama-3.1-8B, Gemma-2-9B).  The component validation results
(\S\ref{sec:validation35b}) suggest that structured tasks (NLI, CDA-7,
claim extraction) degrade substantially below ${\sim}7$B active parameters,
making the 35B MoE backbone the practical minimum for reliable epistemic
fidelity output.  The architectural design of \system\ is backbone-agnostic;
the prompt templates and graph pipeline transfer to any instruction-tuned model.
We expect CR and VC gains to hold qualitatively across stronger backbones,
with improved absolute numbers on NLI-sensitive metrics.

\paragraph{Controversy taxonomy thresholds.}
The ED/PI thresholds defining resolved/emerging/stable/polarized are
design parameters calibrated on the education domain and applied
uniformly.  While per-domain results (Table~\ref{tab:perdomain}) show
consistent behaviour across five domains, optimal thresholds may vary
for domains with fundamentally different disagreement structures
(e.g., legal vs.\ scientific).

\paragraph{NLI noise as dominant error source.}
Error analysis (\S\ref{sec:errors}) identifies prompt-based NLI noise
as the source of $41\%$ of failure cases.  This is an inherent
limitation of zero-shot NLI with instruction-tuned models; a
fine-tuned cross-encoder would likely improve CR by reducing false
contradiction edges in $G$.

\paragraph{Epistemic over-amplification.}
\system\ may over-amplify minority views even when the minority view is
methodologically weak.  Source count display partially mitigates this, but
users in our study still attended to minority views regardless of source count
(see \S\ref{sec:human}): annotators rated minority-view salience as notable
even when those views were supported by a single source.  Stronger visual
uncertainty communication (e.g., proportional-area view displays) remains an
open interface design problem.  We do not recommend deploying \system\ in
contexts where minority-view amplification could cause harm --- for example,
queries about vaccine safety where surfacing methodologically weak dissenting
views may contribute to misinformation, even when those views are labelled with
low confidence and few sources.

\paragraph{CDA-7 resolvability priors.}
The seven resolvability priors ($\rho$ values) are design parameters grounded
in STS literature \citep{longino2002science,solomon2001social} that captures
decades of analysis of how different types of scientific disagreements tend to
resolve.  The ordering (methodological $\rho{=}0.75$ > replication $0.70$ >
statistical $0.60$ > population $0.50$ > operational $0.40$ > temporal $0.35$
> theoretical $0.30$) reflects the widely documented empirical pattern that
methodological and replication disagreements are more often resolved by follow-on
studies than theoretical disagreements, which may persist indefinitely
\citep{longino2002science}.  We acknowledge these are not empirically estimated
from a held-out dataset; Bayesian calibration against a longitudinal corpus of
resolved vs.\ persisting controversies would strengthen the grounding.
Critically, the downstream system behavior is not sensitive to small prior changes:
the confidence output changes class (e.g., medium $\to$ low) only when the
$\rho$ shift crosses the class boundary, and the CDA-7 attribution inter-annotator
agreement ($\kappa{=}0.74$, \S\ref{sec:human}) confirms that the taxonomy classes
themselves are reliably identified regardless of the specific prior values.

\paragraph{Ethics: weaponizing surfaced disagreement.}
CDA-7 attribution and the converging/diverging TDC trajectory provide
mitigation against misuse (e.g., misrepresenting settled science on
vaccines as ``actively debated'') by explicitly marking which
controversies are \emph{resolving} vs.\ \emph{growing} and assigning
\textsc{resolved} queries a \textsc{high} confidence label.  The user
trust calibration experiment (\S\ref{sec:human}) confirms that users
appropriately distinguish stable from polarized controversies when
given \system's structured output.  However, bad-faith actors who
cherry-pick minority-view outputs remain a concern that no technical
mitigation fully addresses.

%% ─────────────────────────────────────────────────────────────────────────────
\section{Conclusion}
\label{sec:conclusion}

Standard single-view scientific RAG can fail not because models miss
evidence, but because they smooth genuinely contested evidence into one
answer.  We formalized this as
\emph{epistemic collapse} and showed analytically and empirically that it is
structurally unavoidable for single-view systems and actively amplified
by topical-relevance retrieval.  We introduced \system, a seven-stage RAG
framework for epistemic fidelity, with a four-metric formal suite (CCS, ED,
PI, EC@K), a CDA-7 causal attribution taxonomy grounded in STS literature,
a four-class controversy taxonomy, temporal epistemic drift tracking, and
\bench\ --- an evaluation benchmark for epistemic fidelity in scientific QA
(250 queries, 5 domains, 75 papers).

Across 250 queries spanning education, biomedicine, economics, earth
sciences, and nutrition, \system\ Full achieves CR$\,{=}\,0.742$
vs.\ vanilla RAG's $0.572$ ($+0.170$; $p{<}0.001$),
$+0.124$ over MADAM-RAG ($p{=}0.003$),
and $+0.189$ over PaperQA2 ($p{<}0.001$).
Dense-encoder VC rises from $0.423$ to $0.847$ ($p{<}0.001$),
empirically corroborating the structural single-view coverage bound.
CCS improves from $0.571$ to $0.518$ ($p{=}0.008$).
Human annotators ($n{=}750$ ratings, Krippendorff's $\alpha{=}0.72$)
rate \system\ Full at $4.3/5$ vs.\ $2.1/5$ for vanilla RAG;
CCS correlates with human epistemic completeness at $\rho{=}{-}0.82$
($p{<}0.001$, $n{=}250$ system--query pairs), confirming the
metric captures genuine information loss.
A user trust calibration experiment shows that \system's structured
confidence labels appropriately modulate trust: users assign lower
overturning probability to stable controversies ($32\%$) than
polarized ones ($58\%$; $p{<}0.001$), while vanilla RAG users show
no such discrimination.

Component validation at 35B scale confirms: NLI F1$\,{=}\,0.933$,
CDA-7 $\kappa{=}0.81$ (exceeding human inter-annotator agreement
of $\kappa{=}0.74$), and 100\% schema compliance.
Error analysis identifies NLI noise as the dominant failure mode
($41\%$ of errors), suggesting that fine-tuned NLI models would
yield further CR gains.  Fast offline mode provides
$25.62{\times}$ latency improvement with no output quality loss.

\bibliography{custom}

%% ─────────────────────────────────────────────────────────────────────────────
\appendix

\section{Baseline Implementation Details}
\label{app:baselines}

All baselines and ablation variants share the same chunked corpus
(3,247 text chunks from 75 papers across five domains) and the same
LLM backbone (\texttt{qwen3.6:35b-a3b} via Ollama, chain-of-thought
disabled, \texttt{think=False}).  Embeddings use \texttt{all-MiniLM-L6-v2}
($d{=}384$) for retrieval and metric computation throughout.

\paragraph{Vanilla RAG.}
Dense retrieval with cosine similarity; top-$k{=}10$ chunks retrieved per
query.  A single synthesis call with the prompt: \textit{``Based on the
following passages, answer the question concisely and accurately:
[passages] Question: [query]''}.  No multi-view output, no disagreement
detection, no confidence structure.

\paragraph{Single-agent.}
Identical to Vanilla RAG but uses only the Precision agent query
formulation (``retrieve the most directly relevant, high-confidence
evidence for [query]'') rather than the raw query string.  Top-$k{=}10$
retrieval, single-view synthesis.  Isolates the effect of the Precision
agent formulation from the full adversarial retrieval ensemble.

\paragraph{MADAM-RAG.}
We use the open-source MADAM-RAG implementation \citep{wang2025conflicting}
adapted to our corpus.  Each query triggers a structured multi-agent debate:
a Proposer agent generates an initial answer; an Opposer agent challenges it;
a Moderator agent produces a final synthesized answer resolving the conflict.
All agents use \texttt{qwen3.6:35b-a3b} as the backbone (replacing the
original GPT-4 backbone for fair comparison).  Top-$k{=}10$ retrieval;
debate structured over three turns per query.  MADAM-RAG produces a
single resolved answer, not a multi-view output.

\paragraph{PaperQA2.}
We use PaperQA2 version 5.x \citep{lala2023paperqa} adapted to our
pre-indexed corpus rather than its native paper-download pipeline.
The same 75 papers were ingested via the PaperQA2 document API;
retrieval uses PaperQA2's default two-stage retrieval (keyword filter
then dense re-ranking).  The LLM backbone was replaced with
\texttt{qwen3.6:35b-a3b} via a compatible API adapter.  PaperQA2
produces citation-grounded single-view answers; its multi-step chain-of-thought
is preserved as designed.

\paragraph{MMR-RAG.}
Maximal Marginal Relevance retrieval \citep{carbonell1998mmr} with
$\lambda{=}0.5$ (equal weight on relevance and diversity) and
top-$k{=}10$ retrieved chunks.  The diversity re-ranking uses cosine
distance between chunk embeddings (\texttt{all-MiniLM-L6-v2}).
Single-view synthesis using the same prompt as Vanilla RAG.
This baseline isolates the effect of retrieval diversity without any
downstream claim-graph reasoning, causal attribution, or multi-view synthesis.

%% ─────────────────────────────────────────────────────────────────────────────
\section{Qualitative Output Examples}
\label{app:examples}

Tables~\ref{tab:qualitative}--\ref{tab:qualitative2} show representative
system outputs on three queries to concretely illustrate epistemic collapse.
Both responses are fluent and locally accurate; the distinction is structural.

\begin{table*}[t]
\centering\small
\setlength{\tabcolsep}{5pt}
\begin{tabular}{@{}p{0.46\textwidth}p{0.50\textwidth}@{}}
\toprule
\textbf{Vanilla RAG} & \textbf{\system\ Full} \\
\midrule
\multicolumn{2}{@{}l}{\textit{Query: Does homework improve academic achievement?}} \\
\midrule
``Research consistently shows that homework improves academic
achievement, with stronger effects for older students (grades 7--12).
Meta-analyses report moderate positive correlations between homework time
and academic performance.  Homework builds self-regulation and reinforces
classroom learning.''

\smallskip
\textit{CCS: 0.38.\ \ VC$_\text{emb}$: 0.33 (1/3 views covered).}
&
\textbf{View 1 [Dominant, 6 sources, confidence: medium]:} Conditional
positive effect.  Homework improves achievement for grades 7--12 ($d{\approx}0.30$);
effect near zero for primary grades.
\textit{CDA-7: methodological (grade-level scope differences).}

\textbf{View 2 [Alternative, 3 sources, confidence: medium]:} Null
meta-analytic effect when controlling for SES and teacher quality.
\textit{CDA-7: statistical + methodological.}

\textbf{View 3 [Minority, 2 sources, confidence: low]:} High homework
loads harm wellbeing with no academic benefit in high-achieving populations.
\textit{CDA-7: operational (achievement vs.\ wellbeing outcomes).}

\smallskip
\textit{Controversy class: Stable.\ \ CCS: 0.52.\ \ VC$_\text{emb}$: 1.00.}
\\
\bottomrule
\end{tabular}
\caption{Homework query: Vanilla RAG collapses to the dominant conditional-benefit
view; \system\ Full recovers all three ground-truth viewpoints with per-view
source counts, CDA-7 causal attribution, and evidence-structured confidence.}
\label{tab:qualitative}
\end{table*}

\begin{table*}[t]
\centering\small
\setlength{\tabcolsep}{5pt}
\begin{tabular}{@{}p{0.22\textwidth}p{0.36\textwidth}p{0.36\textwidth}@{}}
\toprule
\textbf{Query} & \textbf{Vanilla RAG} & \textbf{\system\ Full (summary)} \\
\midrule
\textit{Do statins prevent cardiovascular events in primary prevention?}
&
``Statins are recommended for high-risk primary prevention and reduce
cardiovascular events.'' (1 view; VC$_\text{emb}$: 0.33; CCS: 0.41)
&
View 1 [Dominant]: Benefit confirmed for high-risk groups
\textit{(CDA-7: population)}.
View 2: Marginal benefit with adverse-effect trade-off
\textit{(CDA-7: statistical)}.
View 3: Industry-bias concern undermining reported effect sizes
\textit{(CDA-7: methodological)}.
VC$_\text{emb}$: 1.00; Controversy class: Stable; CCS: 0.54.
\\[4pt]
\textit{Does raising the minimum wage reduce employment?}
&
``Most studies find minimal employment effects from moderate minimum wage
increases.'' (1 view; VC$_\text{emb}$: 0.50; CCS: 0.44)
&
View 1 [Dominant]: Minimal aggregate employment effect
\textit{(CDA-7: methodological)}.
View 2: Significant job losses for low-wage workers in specific sectors
\textit{(CDA-7: population)}.
View 3: Positive employment via demand-side stimulus
\textit{(CDA-7: theoretical)}.
VC$_\text{emb}$: 1.00; Controversy class: Polarized; CCS: 0.50.
\\
\bottomrule
\end{tabular}
\caption{Additional examples in biomedicine and economics.  In both cases
Vanilla RAG returns a single consensus statement that erases 2/3 of
ground-truth viewpoints; \system\ Full provides structured multi-view output
with per-view CDA-7 attribution and controversy-class labels.}
\label{tab:qualitative2}
\end{table*}

\end{document}
